\section{Relativization vs Satisfaction} \label{Ch3:Sec:Satisfaction}

Can we formalize into Set Theory the notion of ``$M \models \vp\of{\bar{x}}$''?

Pick ``symbols'' and formalize the following using $\ZFwoP$:
\begin{itemize}
    \item syntax of first-order logic
    \item definition of satisfaction
    \item proof theory
    \item completeness theorem
\end{itemize}

The way we define what ``formulae'' are internally is using Gödel numbers. Ie, for each formula $\vp$ (\textit{external}), there exists some set $\ulcorner \vp \urcorner$ that is analogous to it. We still write $M \models \vp\of{\bar{x}}$.

It is possible to show that our analogies translate both ways, ie, $\vp\of{\bar{x}}^M$ iff ``$M \models \vp\of{\bar{x}}$'', simply by structural induction on $\vp$.

% Transcribed from the board photographs supplied with the post-lecture pass of 8 Oct 2026; these boards were not typed up during the lecture.

Note, about $M \models \vp\of{\bar{x}}$:
\begin{enumerate}
    \item $\vp$ is not a set. There's a set $\ulcorner \vp \urcorner$ that is analogous to $\vp$. Nevertheless, we just write $M \models \vp\of{\bar{x}}$.
    \item This requires $M \ne \emptyset$ and $\bar{x} \in M^{n + 1}$, where $v_0, \ldots, v_n$ are the free variables of $\vp$. Contrast with $\vp^M\of{\bar{x}}$.
    \item $\vp^M\of{\bar{x}}$ was defined for each $\vp$ and each class $M$. $M \models \vp\of{\bar{x}}$ is defined only for sets $M$:
    \begin{align*}
        \models \; \ssq \; \setst{\parenth{M, c, \bar{x}}}{M \ne \emptyset, \, c \text{ a Gödel number}, \, \bar{x} \in M^{<\omega}}
    \end{align*}
\end{enumerate}

These agree wherever both make sense.

\begin{boxlemma}[Fact (Scheme)]
    Consider any formula $\vp\of{v_0, \ldots, v_n}$. For each $M \ne \emptyset$ and $\bar{x} \in M^{n + 1}$,
    \begin{align*}
        \vp\of{\bar{x}}^M \text{ iff } ``M \models \vp\of{\bar{x}}"
    \end{align*}
\end{boxlemma}

For now, work in $\ZF$. Given a formula $\vp(u)$ and a set $x$, say $\vp(u)$ \textbf{defines} $x$ iff
\begin{align*}
    \forall u \brac{u = x \lr \vp(u)} \tag{$\star$}
\end{align*}
Observe: there is no formula $\delta\of{r, x}$ so that, for every formula $\vp(u)$,
\begin{align*}
    \forall x \brac{\delta\of{\ulcorner \vp(u) \urcorner, x} \lr \underbrace{\vp(u) \text{ defines } x}_{(\star)}}
\end{align*}
Otherwise, the least undefinable ordinal would exist and be definable (because there are only countably many Gödel numbers).
